\scanpage{124}

\[
\left\{\begin{aligned}
\deleted{F}\Phi_{2p}(T)&\inserted{=F_p(-T)\deleted{\unsure{(-1)^{\varphi'(p)}}}}
=1-T+T^2+\ldots+(-1)^{p-1}T^{p-1}\qquad\text{$p$ nb premier impair}\\
\deleted{\Phi}F_{2p}(T)&=\inserted{(-1)^{p'}}\Phi_p(-T)
=\bigl(S_0(T)-S_1(T)+\ldots+(-1)^{p'}S_{p'}(T)\bigr)(-1)^{p'}
\end{aligned}\right.
\]

\begin{center}\rule{1cm}{.4pt}\end{center}

Sur un corps \inserted{$k$ alg clos de} car. quelconque $p$, les racines de $F_n(T)$ sont les racines \uline{primitives} \deleted{\unsure{$n$-ièmes}} \inserted{\deleted{i.e.} $n'$-ièmes} de 1, chacune avec leur multiplicité, qui est

\[
\left\{\begin{aligned}
&\text{multiplicité }1\quad\text{si }(p,n)=1\\
&\text{multiplicité }\varphi(p^r)\quad\text{si }n=p^rn',\quad(n,n')=1,\quad r\in\mathbb N
\end{aligned}\right.
\]

Si $(\zeta_i)\deleted{\unsure{_{1\leq i}}}$ sont les racines \inserted{primitives \deleted{\unsure{distinctes $n'$ de 1|distinctes $n$-ièmes}}} \deleted{\unsure{deux à deux}} $n'$ièmes de 1, on a donc

\[
\Phi_n(T)=\left(\prod(T-\zeta_i)\right)^{\varphi(p^r)}=\Phi_{n'}(T)^{\varphi(p^r)}\quad\text{dans }k[T]
\]

\deleted{\unsure{(il vient donc pour $k$ corps de car. $p$}} on a

\[
\Phi_n(T)=\Phi_{n'}(T)^{\varphi(p^r)}\quad\text{dans }\unsure{\Phi_p(T)|\mathbb F_p(T)}
\]

i.e

\[
\Phi_n\equiv\Phi_{n'}^{\varphi(p^r)}\quad\bmod p
\]

En particulier

\[
\left\{\begin{aligned}
\Phi_{p^r}(T)&\equiv(T-1)^{\varphi(p^r)}\quad\bmod p\\
\Phi_p(T)&\equiv(T-1)^{p-1}\quad\bmod p\qquad -
\end{aligned}\right.
\]

\deleted{\unsure{Les}} Les racines de $F_n(T)$ \deleted{\unsure{dans}} $k$ sont les \unsure{éléments|élémts} de la forme $\zeta+\zeta^{-1}$ ($\zeta$ racine primitive $n'$-ième de 1), chacun avec la multiplicité $\varphi(p^r)\deleted{\unsure{/2}}$ -- \deleted{\unsure{donc}} \inserted{sauf si} $n'=1$ ou 2, où la multiplicité est $\varphi(p^r)/2$.

\[
\left\{\begin{aligned}
F_n(T)&\equiv F_{n'}(T)^{\varphi(p^r)}\quad\bmod p\qquad\text{si }n'\ne1,2\\
F_n(T)&\equiv F_{n'}(T)^{\varphi(p^r)/2}\quad\bmod p\qquad\text{si }n'=1\text{ ou }2
\end{aligned}\right.
\]

en particulier

\[
\left\{\begin{aligned}
F_{p^r}(T)&\equiv\deleted{\unsure{F_p(T)^{\varphi(p)}}}\,(T-2)^{\varphi(p^r)}\quad(p)\\
F_p(T)&\equiv(T-2)
\end{aligned}\right.
\]

% archive sheet number: 123
